\documentclass[11pt]{article}
\usepackage[T1]{fontenc}
\usepackage{amsmath,amssymb,amsthm,mathtools}
\usepackage[a4paper,margin=1in]{geometry}
\usepackage[hidelinks]{hyperref}
\usepackage{microtype}
\usepackage{booktabs}

\newcommand{\F}{\mathbb F}

\title{Boundary-Shift Aggregation of Low-Weight Frontier Polymers}
\author{Research note}
\date{28 July 2026}

\begin{document}
\maketitle

\begin{abstract}
We preserve boundary-shift aggregation in the low-weight polymer relaxation
for Frontier cap pruning.  An exact XOR dynamic program groups all subsets of
visible size-one/two polymers by their active-detector-plus-logical shift
before taking the fractional power.  On the retained surface distance-3
models, this recovers up to 34 orders of magnitude at fixed fractional
exponent.  After optimizing the exponent, however, the partial right-hand
side remains above one at \(K=1024\) and \(K=2048\).  Quotient aggregation
alone is therefore insufficient; compatibility or the exact comparison
spectrum must be retained next.
\end{abstract}

\section{Nested low-weight relaxations}

Fix a cut \(t\), Chernoff order \(\theta\), and let
\(\mathcal P_t^{\le2}\) be the visible size-one/two open-prefix polymers.
For polymer \(\gamma\), write
\[
 w_\gamma=\prod_{j\in\gamma}a_{\theta,j},
 \qquad
 s_\gamma=D_t\mathbf1_\gamma.
\tag{1}
\]
Drop compatibility, but preserve the XOR boundary shift:
\[
 Z_t^{\le2}(\xi)
 =
 \sum_{\substack{
 \mathcal F\subseteq\mathcal P_t^{\le2}\\
 \bigoplus_{\gamma\in\mathcal F}s_\gamma=\xi}}
 \prod_{\gamma\in\mathcal F}w_\gamma.
\tag{2}
\]
For \(0<\rho<1\), define
\[
\begin{aligned}
 G_{t,\rho}^{\le2}
 &=
 \sum_{\xi\ne0}Z_t^{\le2}(\xi)^\rho,\\
 U_{t,\rho}^{\le2}
 &=
 \prod_{\gamma\in\mathcal P_t^{\le2}}(1+w_\gamma^\rho)-1,\\
 E_{t,\rho}^{\le2}
 &=
 \exp\left(\sum_{\gamma\in\mathcal P_t^{\le2}}w_\gamma^\rho\right)-1.
\end{aligned}
\tag{3}
\]
Then
\[
\boxed{
 G_{t,\rho}^{\le2}
 \le U_{t,\rho}^{\le2}
 \le E_{t,\rho}^{\le2}.}
\tag{4}
\]
The first inequality groups nonempty families by shift and uses
\((\sum_i x_i)^\rho\le\sum_i x_i^\rho\), while excluding nonempty
zero-shift families.  The second is \(\log(1+x)\le x\).

With polymers of every size, dropping compatibility from the
bounded-hypergraph pointwise theorem gives
\[
 \omega_t(\xi)\le Z_t(\xi),
\tag{5}
\]
and hence
\[
 D^{\rm cap}
 \le K^{-1/\rho}\sum_t
 \left[\sum_{\xi\ne0}Z_t(\xi)^\rho\right]^{1/\rho}.
\tag{6}
\]
The computed \(Z_t^{\le2}\) is a nonnegative partial contribution, because
adding activity \(w\) at shift \(s\) updates
\[
 Z'(\xi)=Z(\xi)+wZ(\xi+s)\ge Z(\xi).
\tag{7}
\]
It is therefore a lower bound on the right-hand side of (6), not on actual
cap loss.  Compatibility remains unenforced.

\section{Exact XOR dynamic program}

Encode active detector bits followed by logical bits.  Surface d3 has at
most 11 active detectors and one logical bit, hence at most \(2^{12}=4096\)
quotient labels.  Initialize
\[
 Z^{(0)}(0)=1,\qquad Z^{(0)}(\xi\ne0)=0.
\]
For polymer \(m\), activity \(w_m\), and shift \(s_m\), use
\[
\boxed{
 Z^{(m+1)}(\xi)
 =
 Z^{(m)}(\xi)+w_mZ^{(m)}(\xi+s_m).}
\tag{8}
\]
This exactly evaluates (2) in
\(O(|\mathcal P_t^{\le2}|2^{|\Gamma_t|+k})\) time per cut.  Final cut
integration is performed in the log domain.

\section{Surface result}

The deterministic profile uses
\[
 p_{\rm location}=0.001,\qquad
 \alpha=0.8,\qquad
 \theta=5/16,
\]
with \texttt{\detokenize{deadline_reorder}}.  The ratio between the old
exponential expression and the shift-grouped expression, after rooted cut
integration, is
\[
\begin{array}{lrrr}
\toprule
\text{scope}&\rho=0.50&\rho=0.75&\rho=0.99\\
\midrule
\text{memory X}&2.13\,10^{34}&3.03\,10^4&1.215\\
\text{memory Z}&1.18\,10^{28}&7.81\,10^3&1.207\\
\bottomrule
\end{array}
\tag{9}
\]
At \(\rho=0.5\) and \(K=1024\), X falls from
\(1.20\,10^{35}\) to \(5.63\), and Z falls from
\(6.12\,10^{28}\) to \(5.19\).  Aggregation removes enormous overcount,
but the partial expression remains above one.

The artifact also searches
\(\rho=0.05,0.06,\ldots,0.99\):
\[
\begin{array}{rrrrr}
\toprule
K&\text{X minimum}&\rho_X&\text{Z minimum}&\rho_Z\\
\midrule
16&285.558&0.99&269.676&0.99\\
512&8.617&0.99&8.137&0.99\\
1024&4.183&0.85&3.859&0.80\\
2048&1.395&0.45&1.299&0.50\\
4096&0.0715&0.05&0.116&0.05\\
8192&6.82\,10^{-8}&0.05&1.11\,10^{-7}&0.05\\
\bottomrule
\end{array}
\tag{10}
\]
The grid-estimated continuous partial thresholds are
\[
 K_X=2346.45\quad(\rho=0.37),
 \qquad
 K_Z=2317.66\quad(\rho=0.45).
\tag{11}
\]
These are grid diagnostics, not proofs of the continuous optimum.  The
\(K=1024\) conclusion is nevertheless robust: every tested exponent leaves
the size-one/two expression at least \(3.86\), and higher polymers can only
increase it.

At \(K=4096\), the literal cap cannot discard a surface-d3 boundary label
because there are at most 4096 active-plus-logical labels.  The nonzero
fractional expression there only shows that the fractional corollary does
not use the finite-support endpoint sharply.

\section{Conclusion and reproduction}

Boundary aggregation approximately halves the grid-estimated low-weight
threshold: from 5235 to 2346 for X and from 4914 to 2318 for Z.  By itself it
remains insufficient at \(K=1024\) and \(K=2048\).

The subsequent audit is completed in
\texttt{\detokenize{ADAPTIVE_KERNEL_SPECTRUM.tex}}.  Exact compatibility,
shift-specific Chernoff orders, an all-size completed-kernel recurrence, and
exact head trimming give rigorous surface-d3 cap bounds of \(0.2519\) for X
and \(0.2767\) for Z at \(K=1024\).  The result is not FER and does not cover
score-gap pruning or terminal logical ranking.  BB144/Gross remains beyond
the dense exact-spectrum calculation.

\begin{verbatim}
frontier-overlap-profile \
  --backend rotated_surface_d3 --p-location 0.001 \
  --column-order deadline_reorder --score-alpha 0.8 \
  --K-values 16,512,1024,2048,4096,8192 \
  --aggregate-boundary-shifts --max-boundary-bits 16 \
  --out docs/data/overlap_profiles/\
rotated_surface_d3_p0p001_boundary_aggregated.json
\end{verbatim}
The expanded artifact remains deterministic, uses one CPU, and prints
cut-count, elapsed-time, and ETA progress.  The current runtime and SHA256
are recorded in \texttt{\detokenize{WORKLOG.tex}}.  Regression tests compare
(8) with brute subset enumeration, verify (4) cut by cut, and independently
check the later compatibility and kernel recurrences.  No decoder or Monte
Carlo sampling is used.

\end{document}
